\documentclass[11pt]{article}
\usepackage[margin=1in]{geometry}
\usepackage{amsmath,amssymb,amsthm,booktabs,microtype,hyperref}
\hypersetup{hidelinks,pdfauthor={T},pdftitle={Minimum Revalidation Across Uncertain Dependency Histories}}
\newtheorem{theorem}{Theorem}
\title{Minimum Revalidation Across Uncertain Dependency Histories}
\author{T}
\date{}
\begin{document}
\maketitle
\begin{abstract}
Restoring an agent after a compromise requires checking which artifacts
may have inherited suspect state. When incomplete logs admit several
incompatible dependency histories, their edge union can invent paths and
inflate revalidation cost without bound. We characterize when the union
preserves the feasible recovery plans, establish the complexity of the
scenario-wise problem, and give an exact witness-driven solver. An
independent checker can verify optimality using path branches and
rational prices. An alternative encoding shares repeated rebuild
sets. The tests distinguish the size of a proof from the size of
its encoding.
\end{abstract}

\section*{Summary and scope}
\paragraph{Question.} When can a recovery tool replace several
plausible dependency histories with their edge union?
\paragraph{Finding.} The replacement is exact if every union path
contains the eligible vertex set of a path from a listed scenario.
Otherwise the union can overstate minimum recovery
cost by an arbitrarily large factor. A witness search and a separate
certificate checker handle the scenario-wise problem directly.
\paragraph{Evidence and scope.} The equivalence and certificate
claims are proved, with finite exact checks against independent
reachability and optimization oracles. The checker verifies the
admitted histories. Whether those histories cover a real incident
requires provenance evidence outside the optimization.

\section{Revalidation as a recovery action}
After a compromised tool adapter has been isolated, an operator needs
to decide which artifacts to reconstruct before resuming service.
Examples include cached tool outputs, serialized plans, and derived
configuration. A reconstruction may be expensive even when the original
artifact is small because it must start from authorized, independently
validated inputs.

Let $\mathcal G$ be a finite, explicit family of plausible directed
dependency graphs on a common vertex set. An edge represents possible
dependence within the history it belongs to. Let $s$ denote the origin
of suspect state and $t$ denote a target for resumption. Multiple
origins or targets can be represented by a super-source or super-sink.
Eligible vertices form a set $U$ disjoint from $\{s,t\}$. Choosing
$v\in U$ at cost $c(v)>0$ replaces its state with a version independently
revalidated against the authorized specification. The replacement
discards suspect incoming dependencies and preserves downstream access
to the new version. As far as transmission of existing taint is
concerned, this operation removes $v$ from every $s$-to-$t$ path.
Every true dependency history must belong to $\mathcal G$ for a
guarantee about the actual recovery to follow.

Robust cut problems use multiple graph scenarios in several ways.
Kakimura and Schlotter \cite{kakimura} study proximity to
scenario-specific optima. Here the same set of reconstructions
must interrupt all suspect paths in every scenario.

\section{When the union is an exact recovery model}
For a path $P$, let $I(P)=V(P)\cap U$. A revalidation set $R$
works in a graph if and only if it intersects $I(P)$ for every
$s$-to-$t$ path in that graph. Equivalently, the target is unreachable
from the suspect source after removing $R$. Paths with no eligible
internal vertex certify that revalidation of these artifacts alone
cannot authorize resumption.

\begin{theorem}
A cut in the edge union of $\mathcal G$ always works in every
scenario. The converse holds for every $R\subseteq U$ exactly
when, for each union path $Q$, some path $P$ in a single scenario
satisfies $I(P)\subseteq I(Q)$.
\end{theorem}
\begin{proof}
All scenario paths occur in the union, which gives the first
implication. Under the stated condition a scenario-wise cut
intersects $I(P)$ and therefore intersects $I(Q)$.
Conversely, if a union path $Q$ has no such scenario path, take
$R=U\setminus I(Q)$. Each scenario path has an eligible vertex
outside $I(Q)$, so $R$ hits it. The set $R$ misses $Q$.
\end{proof}

Thus the union can be exact even if some paths never occur in a
single scenario, provided every such path contains a sufficient
scenario constraint. Only inclusion-minimal eligible path sets
matter to this test. If the condition fails and $I(Q)$ is nonempty,
giving vertices in $I(Q)$ cost $M>|U|$ and all others cost one
forces the union optimum to exceed the scenario-wise optimum or
be infeasible. If $I(Q)$ is empty, the union is infeasible.

For a concrete family, put the path $s\to r\to t$ in both
scenarios. Add $s\to a$ and $b\to t$ to the first and $a\to b$
to the second. Each scenario has only the path through $r$.
The union also contains $s\to a\to b\to t$. Set $c(r)=1$
and $c(a)=c(b)=M$. The optimum costs are one and $1+M$,
respectively. The gap follows from an incompatible provenance
hypothesis, even though each edge is individually possible.

\section{Computing a revalidation plan}
A single scenario is a weighted vertex cut. Split each eligible
vertex into an entrance and an exit joined by a capacity $c(v)$
edge. Give other edges capacity greater than the total eligible
cost. Any affordable cut uses only the split edges and yields
a revalidation plan by max-flow min-cut \cite{ford}.
If even deleting all eligible vertices fails to disconnect the
target, the problem is infeasible.

For multiple scenarios, start with $R=\varnothing$. If some
graph still connects $s$ to $t$, return one witness path $P$.
Branch on adding each vertex of $I(P)\setminus R$. Every
extension must choose at least one of these vertices, so taking
the least-cost terminal solution is exact. If no path remains,
the current set is a certificate. If the candidate branch set
is empty, the current branch is infeasible. Caching by $R$ is
safe since the witness search depends only on the scenario
family and $R$. With unit costs, at most $\ell$ eligible
vertices per witness and a size cap $k$, search uses at most
$1+\ell+\cdots+\ell^k$ reachability rounds. The parameter bound
is useful when the smallest actionable witness has few eligible
components, or when all witness paths are short.

The general decision problem is NP-complete, even with unit
costs and a single three-edge path per scenario. Given a graph
for vertex cover \cite{karp}, make each graph vertex eligible.
For every undirected edge $\{u,v\}$, choose its orientation
using a common total order and create one scenario with path
$s\to u\to v\to t$. A revalidation set works precisely when
it is a vertex cover. The reduction is polynomial, and a
candidate solution is verified by reachability in each scenario.


\section{Checking optimality without trusting the solver}
A proposed set $R$ of cost $C$ gives an upper bound after
reachability confirms that it blocks the target in every graph.
For positive integer costs, proving optimality then amounts to
ruling out every set of cost at most $C-1$. Integer programming
certificates support independent checks of such bounds \cite{vipr}.
Our certificate uses only operations visible in the recovery model.

At a node, some rebuilds $S$ are already chosen and the remaining
allowance is $b=C-1-\sum_{v\in S}c(v)$. A witness path must
appear in one admitted scenario and avoid $S$. A branch node
lists such a path and supplies one child for each eligible vertex
on it. Every feasible completion must choose at least one of
those vertices. Two kinds of leaf need only a local check.
A negative allowance excludes the branch by cost, and a witness
path with no eligible vertex excludes it by reachability.

A third kind of leaf uses nonnegative rational prices $y_P$ on
a finite list of residual witness paths. For each eligible
vertex $v$, require
\[
\sum_{P:\,v\in I(P)} y_P\le c(v).
\]
Every completion $T$ that hits those paths obeys
\[
\sum_{P}y_P
\le \sum_P y_P\,|I(P)\cap T|
= \sum_{v\in T}\sum_{P:\,v\in I(P)}y_P
\le \sum_{v\in T}c(v).
\]
Therefore a total price strictly greater than $b$ excludes the
branch. The checker verifies the edges of each cited path against
its named graph, so a path fabricated by mixing histories cannot
justify a lower bound.

\begin{theorem}
A feasible $R$ of integer cost $C$ and a valid lower-bound
certificate for allowance $C-1$ prove that $R$ is optimal.
Whenever $R$ is optimal, a finite certificate exists using only
branches, negative-allowance leaves and paths with no eligible
vertex.
\end{theorem}
\begin{proof}
Soundness follows by induction from the local tests and branch
coverage. For completeness, branch on an available witness path
as long as the allowance is nonnegative and its eligible set is
nonempty. Each branch adds a vertex. If a branch within allowance
had no witness path, it would be a feasible plan cheaper than
$R$, contradicting optimality. The finite tree therefore closes
at one of the two elementary leaf types. Rational prices are
optional shortcuts and do not affect completeness.
\end{proof}

For three unit-cost vertices with residual path sets
$\{a,b\},\{a,c\},\{b,c\}$, pricing every path at $1/2$
gives lower bound $3/2$. No unit-cost solution is possible.
Any two vertices suffice. With two disjoint copies of this
triangle, the true cost is four but root path prices cannot
exceed three. Indeed, summing the six unit vertex capacities
counts each path price twice. A branch can still complete the
proof. The two examples distinguish rounding a lower bound
from assuming the fractional problem solves every instance.

An independent checker validated the upper and lower certificates
on 128 seeded graph pairs. It rejected five mutants, including
a fabricated edge, an excessive price, a missing branch, a
false cost prune and an infeasible proposed set. Pricing
reduced the total number of nodes from 330 to 136.
However, the compact JSON representation grew from 6,511 to
7,070 bytes across the fixture. Prices saved bytes in 46 cases,
increased them in 81, and tied in one. A one-node certificate
can carry several paths, so node count alone is not a useful
storage guarantee. An operator with a strict evidence budget
can compare the sizes of two verified certificates before
choosing which to store. The verifier checks the stated model.
It cannot establish that the admitted dependency histories
cover what happened in the incident.

\subsection{Sharing proofs for the same rebuild set}
For fixed histories, costs and starting allowance, a partial proof
state is determined by the chosen set $S$. The allowance is the
starting allowance minus $c(S)$, and the residual paths are those
avoiding $S$. A verifier can therefore read one record per set $S$.
At a branch with witness path $P$, the required child records are
precisely $S\cup\{v\}$ for $v\in I(P)$. Each child adds a vertex,
so the graph is acyclic. Induction in decreasing order of $|S|$
proves the same lower bound as the tree. For completeness, choose
one valid subtree for each repeated state in a complete tree and
retain only the records reachable from the empty set. Each shared
record still has all the children its witness requires.

There is an exact separation for $k$ disjoint triangles, with one
scenario per edge and budget $2k-1$. Process the triangles in order.
Within each triangle, branch first on $\{a,b\}$, then on the edge
opposite the vertex just chosen. The tree has four occurrences of
a two-vertex completion, while there are only three distinct pairs.
Writing $T_k$ for the number of tree nodes and $D_k$ for the number
of distinct states gives $T_0=D_0=1$ and
\[
T_k=3+4T_{k-1}=2\cdot 4^k-1,\qquad
D_k=3+3D_{k-1}=\frac{5\cdot 3^k-3}{2}.
\]
The construction uses unit costs and a separate path $s,u,v,t$
for each triangle edge $\{u,v\}$. Its optimum is $2k$ because
any plan needs two vertices in each triangle.

In compact JSON, the five-triangle example used 2,047 tree nodes
and 606 distinct states. A tree with a shared path dictionary
used 9,798 bytes. Adding state sharing reduced it to 6,257 bytes,
a reduction of 36.1\%. The DAG was larger for one and two
triangles. On the 128 seeded weighted graph pairs, it shared no
states and used 5,069 bytes in total against 4,125 for the
path-dictionary tree. The two encodings therefore favor different
fixtures. The checker rejected seven negative controls, including
an omitted child, a forged path, and use of the proof at a higher
allowance. It also accepted a direct source-to-target witness
with no eligible vertex.

\section{Exact checks and use in incident response}
The checker compares witness search with subset enumeration
using ordered transitive closure. They agree for both the
scenario-wise and union optima on all 192 seeded pairs. The
union costs strictly more on 69 pairs, with a summed gap of
373 cost units and a maximum gap of 19. These are properties
of the fixed synthetic design, not observed operational rates.
The checker also verifies all 1,024 simple graphs on five
labeled vertices against an independent vertex-cover oracle.
In the explicit gap family it recovers costs one and 501 at
$M=500$.
A separate check enumerates all 1,024 ordered pairs of forward
four-vertex graphs without a direct source-to-target edge.
Path enumeration tests the inclusion condition, while reachability
tests feasible sets for all four choices of $R$. Both methods
identify the same 252 pairs for which the feasible sets differ.

In practice, bind the plan to the set of dependency histories
admitted by the evidence. If an operator learns a new possible
history, run the witness search before reusing the old certificate.
Retain a conservative scenario when a missing log could hide
a path. The opportunity is to remove impossible joint histories
using trusted evidence, which can save full rebuilds. The danger
is excluding a genuine history or treating a contaminated
revalidation tool as independent. A prospective exercise should
blind the solver to a known ground-truth log, inject both missing
and mutually exclusive edges, and measure unnecessary rebuilds
and remaining reachable taint. Use runtime isolation until the
chosen reconstruction and the resumption authorization are checked.

\paragraph{Review status.} Independent manuscript review is pending.

\begin{thebibliography}{9}
\bibitem{ford} L. R. Ford Jr.\ and D. R. Fulkerson.
\emph{Maximal Flow Through a Network}.
Canadian Journal of Mathematics 8, 1956.
\url{https://www.cs.yale.edu/homes/lans/readings/routing/ford-max_flow-1956.pdf}.
\bibitem{karp} Richard M. Karp.
\emph{Reducibility Among Combinatorial Problems}.
Complexity of Computer Computations, 1972.
\url{https://www.cs.umd.edu/~gasarch/BLOGPAPERS/Karp.pdf}.
\bibitem{kakimura} Naonori Kakimura and Ildikó Schlotter.
\emph{Parameterized Complexity of Submodular Minimization under Uncertainty}.
2024. \url{https://arxiv.org/abs/2404.07516}.
\bibitem{vipr} Kevin K. H. Cheung, Ambros Gleixner, and
Daniel E. Steffy. \emph{Verifying Integer Programming Results}.
2017. \url{https://arxiv.org/abs/1611.08832}.
\end{thebibliography}
\end{document}
